\chapter*{致谢}

论文付梓之际，回首在华东师范大学上海国际首席技术官学院的学习历程，心中充满感激。

首先，衷心感谢我的导师XXX教授。从论文选题、研究框架搭建到DMAIC各阶段的逻辑梳理，导师始终以严谨的治学态度给予悉心指导。每当我在SERVQUAL维度扩展与AHP权重确定等关键节点陷入困惑时，导师总能以寥寥数语直指要害，让我豁然开朗。导师对服务质量领域的深刻洞见和对学术规范的严格要求，是我完成这篇论文最坚实的支撑。

感谢华东师范大学上海国际首席技术官学院各位授课老师。两年间系统学习了质量管理和数据分析工具，正是"DMAIC方法论"课程讲授启发了本文的核心思路。感谢MEM中心的各位老师在学习事务上的细致安排，让在职学习得以顺畅推进。

感谢P公司领导和同事的大力支持。特别感谢风险管理部、信贷审批部和客户服务部的配合，使486份客户问卷得以顺利发放与回收；感谢五位行业专家在AHP层次分析中投入的专业判断，他们的实务经验为研究注入了不可替代的实践价值。

感谢论文评审专家和答辩委员会各位老师，你们提出的宝贵意见让论文得以不断完善。

最后，感谢我的家人。在职攻读学位的日子里，无数个周末和夜晚被论文占据，是你们的理解与包容让我能够心无旁骛地完成这篇研究。

本论文虽已尽力，但囿于学识与时间，不足之处在所难免。文中的疏漏与偏颇，概由本人负责。
